%%═══════════════════════════════════════════════════════════════════
%% Lecture 01 — Introduction: What Is Particle Physics?
%%              Classical Field Theory and Its Lagrangian Formulation
%% 77609 — Introduction to Elementary Particles
%% Date:
%% Lecturer: Prof. Yonit Hochberg
%%═══════════════════════════════════════════════════════════════════

\section{Lecture 01 — Introduction: Particle Physics, Classical Field
  Theory, and the Limits of Non-Relativistic Quantum
  Mechanics}\label{sec:lec01}
\begin{center}
\begin{tabular}{ll}
  \textbf{Date:}\quad 22/03/2026   & \\
\end{tabular}
\end{center}
\hrule\vspace{1.5em}
%%──────────────────────────────────────────────────────────────────
\subsection*{Overview}

This lecture opens the course by framing the central question of
particle physics — \emph{what is the Lagrangian of nature?} — and
presenting the Standard Model as the current, symmetry-based answer.
Hochberg argues that non-relativistic quantum mechanics is fundamentally
insufficient because it cannot describe the creation and annihilation of
particles; the correct framework is relativistic quantum field theory
(QFT), valid up to the energies probed so far (a few TeV).
The lecture then introduces the building block that makes this
transition possible: the \emph{classical field}, a quantity depending on
both space and time.  By reviewing the Lagrangian formalism of classical
mechanics and generalizing the role of the generalized coordinate from a
particle position $x(t)$ to a field $\phi(\mathbf{x}, t)$, Hochberg
sketches the logical path from classical mechanics $\to$ classical field
theory $\to$ quantum field theory that the rest of the course will
follow.

%%──────────────────────────────────────────────────────────────────
\subsection*{Definitions}

\dfn{Particle physics / high-energy physics}{The branch of physics that asks: \emph{what are the fundamental laws of
nature?}  The working assumption is that nature obeys the
\emph{principle of minimal action}, so the question reduces to: what
is the Lagrangian (or Hamiltonian) of nature?}

\dfn{Standard Model (SM)}{The current best answer to the above question: a relativistic quantum
field theory built on axioms and symmetry principles that describes all
known elementary particles and three of the four fundamental forces.}

\dfn{Field}{A \emph{field} is a quantity $\phi$ that depends on \emph{both} space
and time:
\begin{equation}
  \phi = \phi(\mathbf{x}, t).
  \label{eq:lec01:field}
\end{equation}
A field may be:
\begin{itemize}
  \item a \emph{scalar field} — a single number at each spacetime point
    (e.g.\ temperature $T(\mathbf{x},t)$, density $\rho(\mathbf{x},t)$);
  \item a \emph{vector field} — a vector at each spacetime point
    (e.g.\ wind velocity $\mathbf{W}(\mathbf{x},t)$, magnetic field
    $\mathbf{B}(\mathbf{x},t)$);
  \item a \emph{tensor field} — a tensor at each spacetime point
    (examples deferred to relativistic notation).
\end{itemize}}

\dfn{Vector field (precise)}{A collection of component fields that \emph{transforms as a vector
under spatial rotations}.  While one may formally decompose
$\mathbf{B} = (B_x, B_y, B_z)$ into three independent scalar
functions, it is the rotational transformation property
\begin{equation}
  B_i \;\xrightarrow{\text{rotation}}\; R_{ij}\,B_j
  \label{eq:lec01:vector-rotation}
\end{equation}
that promotes them to a single vector field.}

\dfn{UV cutoff}{The energy scale (equivalently, the shortest length scale) at which a
given field description ceases to be valid.  For the electromagnetic
field and QED, the UV cutoff is believed to be at or below the
\emph{Planck scale}.  For macroscopic fields (temperature, density),
the UV cutoff is set by the microscopic structure (e.g.\ interatomic
spacing).}

\dfn{IR cutoff}{The complementary scale: the \emph{largest} distance (lowest energy)
at which a given field description remains valid.}

\dfn{Classical field theory}{A dynamical theory whose \emph{generalized coordinates are fields}
$\phi(\mathbf{x},t)$ rather than particle positions $x_i(t)$.  Solving
a classical field theory means finding $\phi(\mathbf{x},t)$ for all
$\mathbf{x}$ and $t$, given initial data — a function rather than a
finite set of numbers.}

%%──────────────────────────────────────────────────────────────────
\subsection*{Key Results}

\thm{Lagrangian of classical mechanics: standard form}{In fundamental physics, the Lagrangian for a system with generalized
coordinates $\{x_i\}$ takes the form
\begin{equation}
  L = L(x_i,\,\dot{x}_i),
  \label{eq:lec01:Lcm}
\end{equation}
with \emph{no} explicit time dependence (to ensure energy conservation)
and dependence on \emph{at most} first time-derivatives (standard
assumption).  The equation of motion follows from the
\emph{Euler–Lagrange equation}:
\begin{equation}
  \frac{d}{dt}\frac{\partial L}{\partial \dot{x}_i}
  - \frac{\partial L}{\partial x_i} = 0.
  \label{eq:lec01:EL}
\end{equation}
Each degree of freedom $x_i$ requires \emph{two} initial conditions
(e.g.\ $x_i(t_0)$ and $\dot{x}_i(t_0)$).}

The central boxed statement of the lecture — the key conceptual shift:
\begin{equation}
  \boxed{
    \underbrace{x_i(t)}_{\text{mechanics}}
    \;\longrightarrow\;
    \underbrace{\phi(\mathbf{x},t)}_{\text{field theory}}
    \qquad
    \text{(generalized coordinate becomes the field)}
  }
  \label{eq:lec01:shift}
\end{equation}

\thm{Why non-relativistic QM is insufficient}{Non-relativistic quantum mechanics conserves particle number: the
normalization $\int|\psi|^2\,d^3x = 1$ encodes the assumption that the
particle is \emph{always present}.  In reality, particles are created
and annihilated (e.g.\ $e^+e^- \to \gamma\gamma$, nuclear decays,
collider events).  Describing such processes requires a framework where
particle number is not conserved — quantum field theory.}

\thm{Scope of QFT}{Relativistic quantum field theory is a good description of nature at
all energy scales probed so far, up to $\mathcal{O}(\text{few TeV})$.
Beyond this scale (towards the Planck scale), the framework may break
down, but no experimental violation has been found.}

%%──────────────────────────────────────────────────────────────────
\subsection*{Derivations}

%% DERIVATION 1: Euler-Lagrange equation in classical mechanics
\subsubsection*{Review: principle of minimal action and the
Euler–Lagrange equation}

The \emph{action} is defined as
\begin{equation}
  S[x_i] = \int_{t_1}^{t_2} L(x_i, \dot{x}_i)\,dt.
  \label{eq:lec01:action}
\end{equation}
The \emph{principle of minimal action} states that the physical
trajectory $x_i(t)$ is the one that makes $S$ stationary,
$\delta S = 0$.  Varying \eqref{eq:lec01:action}:
\begin{align}
  \delta S
  &= \int_{t_1}^{t_2}
     \left(
       \frac{\partial L}{\partial x_i}\,\delta x_i
       + \frac{\partial L}{\partial \dot{x}_i}\,\delta \dot{x}_i
     \right) dt \notag\\
  &= \int_{t_1}^{t_2}
     \left(
       \frac{\partial L}{\partial x_i}
       - \frac{d}{dt}\frac{\partial L}{\partial \dot{x}_i}
     \right)\delta x_i\,dt
     + \left[\frac{\partial L}{\partial \dot{x}_i}\,\delta x_i
       \right]_{t_1}^{t_2}.
\label{eq:lec01:deltaS}
\end{align}
With fixed endpoints $\delta x_i(t_1)=\delta x_i(t_2)=0$, the
boundary term vanishes.  Since $\delta x_i$ is arbitrary, $\delta S=0$
requires \eqref{eq:lec01:EL}.

\noindent\textit{Physical meaning of no explicit time dependence.}
If $L$ has no explicit $t$, then by Noether's theorem the associated
conserved charge is the \emph{Hamiltonian} (energy).  Hence, the
assumption $\partial L/\partial t = 0$ is the mathematical statement
that energy is conserved.

%% DERIVATION 2: Effect of symmetry on Lagrangian — worked examples
\subsubsection*{Symmetry constraints on the Lagrangian: examples}

Hochberg illustrated how symmetry requirements constrain the allowed
form of the Lagrangian.  Take
$L = \tfrac{1}{2}m\dot{x}^2 - V(x)$ in one spatial dimension.

\noindent\textbf{Example 1: $\mathbb{Z}_2$ symmetry $x\to -x$.}
Requiring invariance of $L$ under $x\to -x$ means $V(-x)=V(x)$,
i.e.\ $V$ must be an \emph{even} function of $x$ (equivalently, a
function of $|x|$).

\noindent\textbf{Example 2: $SO(3)$ rotational invariance in 3D.}
Promoting $x\to\mathbf{x}\in\mathbb{R}^3$ and requiring invariance
under $\mathbf{x}\to R\mathbf{x}$ for all $R\in SO(3)$ forces $V$ to
depend only on $|\mathbf{x}|$:
\begin{equation}
  V = V(|\mathbf{x}|) \quad \text{(central potential).}
  \label{eq:lec01:central-potential}
\end{equation}
The general lesson: \emph{symmetry principles constrain and
ultimately determine the allowed form of the Lagrangian.}  This will
be the central strategy for model building later in the course.

%%──────────────────────────────────────────────────────────────────
\subsection*{Examples}

\ex{Familiar fields from everyday life}{\begin{itemize}
  \item \textbf{Temperature} $T(\mathbf{x},t)$: scalar field.
    Valid only at scales larger than the interatomic spacing (UV
    cutoff $\sim$ \AA).
  \item \textbf{Wind velocity} $\mathbf{W}(\mathbf{x},t)$: vector
    field.
  \item \textbf{Population density} $\rho(\mathbf{x},t)$: scalar
    field.  Valid only at scales larger than the typical size of a
    person.
  \item \textbf{Magnetic field} $\mathbf{B}(\mathbf{x},t)$: vector
    field.  Its three components $(B_x, B_y, B_z)$ transform as a
    vector under rotations, which is why we bundle them together
    rather than treating them as three independent scalars.
\end{itemize}}

\ex{Solving mechanics vs.\ solving field theory}{In \emph{classical mechanics} with one degree of freedom, solving the
problem means finding $x(t)$ — a function of one argument — given two
initial conditions $x(t_0)$ and $\dot{x}(t_0)$.

In \emph{classical field theory} in $D$ spatial dimensions, solving
the problem means finding $\phi(\mathbf{x},t)$ — a function of $D+1$
arguments.  The ``initial condition'' is now an \emph{entire function}
$\phi(\mathbf{x},t_0)$ (and its time derivative), i.e.\ infinitely
many numbers.}

%%──────────────────────────────────────────────────────────────────
\subsection*{Connections}

\begin{itemize}
  \item \textbf{Analytical Mechanics.}  Hochberg explicitly stated
    that the course assumes knowledge of analytical mechanics at the
    level of the Lagrangian formalism, Euler–Lagrange equations, and
    the principle of minimal action.  Classical field theory is
    presented as a direct generalization of that framework.
  \item \textbf{Quantum Mechanics (two semesters).}  The course also
    assumes two semesters of QM.  The limitations of non-relativistic
    QM (fixed particle number, non-covariant Schrödinger equation)
    motivate the move to QFT.
  \item \textbf{Electromagnetism.}  Basic knowledge of E\&M is
    assumed.  The electromagnetic field $(\mathbf{E}, \mathbf{B})$
    is used as the canonical example of a physical vector field.
  \item \textbf{Quantum Field Theory.}  This course builds a
    ``baby version'' of QFT: classical field theory $\to$ second
    quantization $\to$ perturbation theory $\to$ Feynman diagrams.
    Students are not assumed to have taken QFT previously.
  \item \textbf{Standard Model.}  The SM is the long-range target.
    The course works towards it by building QED, QCD, and spontaneous
    symmetry breaking from first principles.
  \item \textbf{Grossman \& Nir (draft/book).}  From roughly weeks
    5–6 onward, the course follows a draft by Yossi Nir and Yuval
    Grossman (now also a published book), available on Moodle as a
    PDF for internal course use only.
\end{itemize}

%%──────────────────────────────────────────────────────────────────
\subsection*{To Remember}

\begin{itemize}
  \item The central question of particle physics is: \emph{what is
    the Lagrangian of nature?}  The answer, as of now, is the Standard
    Model.
  \item Nature obeys the \emph{principle of minimal action}; all
    dynamics follows from $\delta S = 0$ with
    $S = \int L\,dt$.
  \item Non-relativistic QM cannot describe particle creation and
    annihilation.  The correct framework is relativistic QFT.
  \item In QFT, particles are \emph{excitations of fields}, not
    fundamental objects in their own right.
  \item A field $\phi(\mathbf{x},t)$ is the generalized coordinate of
    field theory; it replaces the particle position $x(t)$ of mechanics.
  \item Every field description has a \textbf{UV cutoff} (shortest
    scale of validity) and an \textbf{IR cutoff} (longest scale).  For
    QED, the UV cutoff is at or below the Planck scale.
  \item Symmetry requirements constrain the form of the Lagrangian:
    this will be the main tool for model building.
\end{itemize}

%%──────────────────────────────────────────────────────────────────
\subsection*{Exam / HW Hints}

\begin{itemize}
  \item \textbf{HW is mandatory for exam eligibility}: all sets must
    be submitted; at most one set may be skipped.  If an extension is
    needed, email Hochberg in advance — she is flexible if asked.
  \item The \textbf{first HW set} will include: (a) deriving
    additional constraints on Lagrangians from symmetry requirements
    (parity, rotational invariance, etc.), extending the examples
    worked in class; (b) elements of the Lagrangian construction for
    classical mechanics.  Hochberg explicitly pointed to the
    in-lecture examples ($x\to -x$ symmetry, rotational invariance)
    as representative of what the homework will ask.
  \item Hochberg explicitly warned against submitting AI-generated
    solutions without understanding them; the exam requires genuine
    mastery of the material.
  \item \textbf{Next topic}: relativistic notation and natural units.
    A COVID-era Panopto recording (2 hours) may be posted if the
    Thursday session cannot be held live.
\end{itemize}

%%──────────────────────────────────────────────────────────────────
\subsection*{Further Reading}

No specific textbook was cited in this lecture.  The following are recommended:

\begin{itemize}
  \item \textbf{Griffiths, \textit{Introduction to Elementary Particles} (2nd ed.), Chapter~1}
    — ``Historical Introduction to the Elementary Particles'': the cast of
    characters (quarks, leptons, bosons), the four forces, and the logic behind
    the Standard Model.  Ideal companion to this introductory lecture.

  \item \textbf{Griffiths, Chapter~6.1} — Brief derivation of Fermi's Golden Rule
    in the particle-physics context, motivating the $(2\pi)^4\delta^{(4)}$ phase-space
    factor that appears throughout the course.

  \item \textbf{Halzen \& Martin, \textit{Quarks and Leptons}, Chapter~1}
    — ``Quarks and Leptons'': a concise survey of the known particles and their
    quantum numbers ($I$, $J^{PC}$, colour).  Good reading alongside the SM overview
    from this lecture.

  \item \textbf{Tong, \textit{Lectures on Particle Physics}, Chapter~1}
    (freely available at
    \href{https://www.damtp.cam.ac.uk/user/tong/particle.html}{\texttt{damtp.cam.ac.uk/user/tong/particle.html}}):
    a streamlined account of what a relativistic quantum field theory is and why
    particles arise from fields.

  \item \textbf{PDG Summary Tables} —
    \href{https://pdg.lbl.gov/2023/tables/contents_tables.html}{\texttt{pdg.lbl.gov}}:
    the definitive up-to-date listing of all known particles, masses, widths, and
    quantum numbers.  Useful to skim even at this early stage to get a concrete
    sense of what the Standard Model contains.
\end{itemize}

%%──────────────────────────────────────────────────────────────────
